% GENERATED FILE. Edit canonical lessons, then run npm run generate:books.
% canonical-source-hash: fnv1a64:e6deaf9a2626b962
% canonical-lessons: MR-C381-joparyant, MR-C381-toparyant, MR-C381-basarte, MR-C381-jenekarun, MR-C381-janu, MR-R381-review-374-381

\chapter{As Long As}
\label{ch:mr-as-long-as}
\hlchaptermodality{381}

\begin{chapteropening}
\textbf{By the end of this chapter:} \emph{I can say until, as long as, provided that and so that.}

Complete the last lesson of chapter 381: I can say until, as long as, provided that and so that.
\end{chapteropening}

\section[\texorpdfstring{\mr{जोपर्यंत} (joparyant)}{जोपर्यंत (joparyant)}]{\mr{जोपर्यंत} (joparyant) — as long as; until}

\label{lesson:MR-C381-joparyant}



\begin{quote}
Before the new one: say the Marathi for to impose, then the Marathi for a provision.
\end{quote}

\subsection*{You'll want to know: \mr{जोपर्यंत}}
\textbf{\mr{जोपर्यंत}} — \emph{joparyant} — \textquotedblleft{}as long as; until\textquotedblright{}.

\begin{quote}
\mr{जोपर्यंत} \mr{तू} \mr{येत} \mr{नाहीस}, \mr{मी} \mr{जेवणार} \mr{नाही}.
\end{quote}

\begin{quote}
\emph{joparyant tū yet nāhīs, mī jevṇār nāhī.}
\end{quote}

\begin{quote}
I won't eat until you come.
\end{quote}

With a negative it means until: \mr{जोपर्यंत} … \mr{नाही}; a \mr{तो}- clause often answers it.

\begin{cousinweb}
\textbf{\mr{जोपर्यंत}} (\emph{joparyant}) is built from \mr{जो}, \textquotedblleft{}which\textquotedblright{}, and -\mr{पर्यंत}, \textquotedblleft{}up to\textquotedblright{}, from Sanskrit \emph{paryanta}, \textquotedblleft{}limit, edge\textquotedblright{} (Monier-Williams): up to the point which.
\end{cousinweb}

\begin{grammarlens}[title={What you've built}]
One more word for saying what happens if.
\end{grammarlens}

\subsection*{Your turn}
\begin{itemize}
\raggedright
  \item \emph{Say it:} \emph{joparyant}
  \item \emph{Say it:} \emph{joparyant}, once more
  \item \emph{Say it:} \emph{joparyant tū yet nāhīs, mī jevṇār nāhī.} — I won't eat until you come.
  \item \emph{Recall:} say the Marathi for to impose, then the Marathi for a provision, then say \emph{joparyant} again
\end{itemize}

\begin{tcolorbox}[breakable,colback=teal!4,colframe=teal!35!black,title={Before you move on}]
What does \mr{जोपर्यंत} mean? (\textquotedblleft{}As long as; until\textquotedblright{}.) Say it once more.
\end{tcolorbox}
\section[\texorpdfstring{\mr{तोपर्यंत} (toparyant)}{तोपर्यंत (toparyant)}]{\mr{तोपर्यंत} (toparyant) — until then; by then}

\label{lesson:MR-C381-toparyant}



\begin{quote}
Before the new one: say the Marathi for a provision, then the Marathi for as long as.
\end{quote}

\subsection*{You'll want to know: \mr{तोपर्यंत}}
\textbf{\mr{तोपर्यंत}} — \emph{toparyant} — \textquotedblleft{}until then; by then\textquotedblright{}.

\begin{quote}
\mr{बस} \mr{सात} \mr{वाजता} \mr{येईल}, \mr{तोपर्यंत} \mr{आपण} \mr{चहा} \mr{घेऊ}.
\end{quote}

\begin{quote}
\emph{bas sāt vājtā yeīl, toparyant āpaṇ cahā gheū.}
\end{quote}

\begin{quote}
The bus will come at seven; until then, let us have tea.
\end{quote}

It answers \mr{जोपर्यंत}: \mr{जोपर्यंत} … \mr{तोपर्यंत} ….

\begin{cousinweb}
\textbf{\mr{तोपर्यंत}} (\emph{toparyant}) is built from \mr{तो}, \textquotedblleft{}that\textquotedblright{}, and -\mr{पर्यंत}, \textquotedblleft{}up to\textquotedblright{}, from Sanskrit \emph{paryanta}, \textquotedblleft{}limit, edge\textquotedblright{} (Monier-Williams): up to that point.
\end{cousinweb}

\begin{grammarlens}[title={What you've built}]
One more word for saying what happens if.
\end{grammarlens}

\subsection*{Your turn}
\begin{itemize}
\raggedright
  \item \emph{Say it:} \emph{toparyant}
  \item \emph{Say it:} \emph{toparyant}, once more
  \item \emph{Say it:} \emph{bas sāt vājtā yeīl, toparyant āpaṇ cahā gheū.} — The bus will come at seven; until then, let us have tea.
  \item \emph{Recall:} say the Marathi for a provision, then the Marathi for as long as, then say \emph{toparyant} again
\end{itemize}

\begin{tcolorbox}[breakable,colback=teal!4,colframe=teal!35!black,title={Before you move on}]
What does \mr{तोपर्यंत} mean? (\textquotedblleft{}Until then; by then\textquotedblright{}.) Say it once more.
\end{tcolorbox}
\section[\texorpdfstring{\mr{बशर्ते} (baśarte)}{बशर्ते (baśarte)}]{\mr{बशर्ते} (baśarte) — provided that}

\label{lesson:MR-C381-basarte}



\begin{quote}
Before the new one: say the Marathi for as long as, then the Marathi for until then.
\end{quote}

\subsection*{You'll want to know: \mr{बशर्ते}}
\textbf{\mr{बशर्ते}} — \emph{baśarte} — \textquotedblleft{}provided that\textquotedblright{}.

\begin{quote}
\mr{मी} \mr{मदत} \mr{करीन}, \mr{बशर्ते} \mr{तुम्ही} \mr{वेळेवर} \mr{याल}.
\end{quote}

\begin{quote}
\emph{mī madat karīn, baśarte tumhī veḷevar yāl.}
\end{quote}

\begin{quote}
I will help, provided that you come on time.
\end{quote}

Formal; it stands before the condition.

\begin{cousinweb}
\textbf{\mr{बशर्ते}} (\emph{baśarte}) is a loan from Persian \emph{ba-sharṭ}, \textquotedblleft{}on condition\textquotedblright{}, usually traced to Arabic \emph{sharṭ}, \textquotedblleft{}condition\textquotedblright{}.
\end{cousinweb}

\begin{grammarlens}[title={What you've built}]
One more word for saying what happens if.
\end{grammarlens}

\subsection*{Your turn}
\begin{itemize}
\raggedright
  \item \emph{Say it:} \emph{baśarte}
  \item \emph{Say it:} \emph{baśarte}, once more
  \item \emph{Say it:} \emph{mī madat karīn, baśarte tumhī veḷevar yāl.} — I will help, provided that you come on time.
  \item \emph{Recall:} say the Marathi for as long as, then the Marathi for until then, then say \emph{baśarte} again
\end{itemize}

\begin{tcolorbox}[breakable,colback=teal!4,colframe=teal!35!black,title={Before you move on}]
What does \mr{बशर्ते} mean? (\textquotedblleft{}Provided that\textquotedblright{}.) Say it once more.
\end{tcolorbox}
\section[\texorpdfstring{\mr{जेणेकरून} (jeṇekarūn)}{जेणेकरून (jeṇekarūn)}]{\mr{जेणेकरून} (jeṇekarūn) — so that}

\label{lesson:MR-C381-jenekarun}



\begin{quote}
Before the new one: say the Marathi for until then, then the Marathi for provided that.
\end{quote}

\subsection*{You'll want to know: \mr{जेणेकरून}}
\textbf{\mr{जेणेकरून}} — \emph{jeṇekarūn} — \textquotedblleft{}so that\textquotedblright{}.

\begin{quote}
\mr{मी} \mr{लवकर} \mr{निघालो}, \mr{जेणेकरून} \mr{गाडी} \mr{चुकू} \mr{नये}.
\end{quote}

\begin{quote}
\emph{mī lavkar nighālo, jeṇekarūn gāḍī cukū naye.}
\end{quote}

\begin{quote}
I left early so that I wouldn't miss the train.
\end{quote}

Followed by a should-form or a future: \mr{जेणेकरून} … \mr{होईल}.

\begin{cousinweb}
\textbf{\mr{जेणेकरून}} (\emph{jeṇekarūn}) is built from \mr{जेणे}, \textquotedblleft{}by which\textquotedblright{}, an old instrumental form of the relative \mr{जो}, \textquotedblleft{}which\textquotedblright{}, and \mr{करून}, \textquotedblleft{}having done\textquotedblright{}: done in such a way that. Sanskrit had the same pair, the relative \emph{ya-} and its instrumental \emph{yena} (CDIAL 10391).
\end{cousinweb}

\begin{grammarlens}[title={What you've built}]
One more word for saying what happens if.
\end{grammarlens}

\subsection*{Your turn}
\begin{itemize}
\raggedright
  \item \emph{Say it:} \emph{jeṇekarūn}
  \item \emph{Say it:} \emph{jeṇekarūn}, once more
  \item \emph{Say it:} \emph{mī lavkar nighālo, jeṇekarūn gāḍī cukū naye.} — I left early so that I wouldn't miss the train.
  \item \emph{Recall:} say the Marathi for until then, then the Marathi for provided that, then say \emph{jeṇekarūn} again
\end{itemize}

\begin{tcolorbox}[breakable,colback=teal!4,colframe=teal!35!black,title={Before you move on}]
What does \mr{जेणेकरून} mean? (\textquotedblleft{}So that\textquotedblright{}.) Say it once more.
\end{tcolorbox}
\section[\texorpdfstring{\mr{जणू} (jaṇū)}{जणू (jaṇū)}]{\mr{जणू} (jaṇū) — as if}

\label{lesson:MR-C381-janu}



\begin{quote}
Before the new one: say the Marathi for provided that, then the Marathi for so that.
\end{quote}

\subsection*{You'll want to know: \mr{जणू}}
\textbf{\mr{जणू}} — \emph{jaṇū} — \textquotedblleft{}as if\textquotedblright{}.

\begin{quote}
\mr{तो} \mr{असा} \mr{वागला}, \mr{जणू} \mr{काही} \mr{घडलंच} \mr{नाही}.
\end{quote}

\begin{quote}
\emph{to asā vāglā, jaṇū kāhī ghaḍlañc nāhī.}
\end{quote}

\begin{quote}
He behaved as if nothing had happened.
\end{quote}

\mr{जणू} \mr{काही}, as though; often after \mr{असा}: \mr{असा} … \mr{जणू}.

\begin{cousinweb}
\textbf{\mr{जणू}} (\emph{jaṇū}) is a Marathi word whose origin is not settled.
\end{cousinweb}

\begin{grammarlens}[title={What you've built}]
One more word for saying what happens if.
\end{grammarlens}

\subsection*{Your turn}
\begin{itemize}
\raggedright
  \item \emph{Say it:} \emph{jaṇū}
  \item \emph{Say it:} \emph{jaṇū}, once more
  \item \emph{Say it:} \emph{to asā vāglā, jaṇū kāhī ghaḍlañc nāhī.} — He behaved as if nothing had happened.
  \item \emph{Recall:} say the Marathi for provided that, then the Marathi for so that, then say \emph{jaṇū} again
\end{itemize}

\begin{tcolorbox}[breakable,colback=teal!4,colframe=teal!35!black,title={Before you move on}]
What does \mr{जणू} mean? (\textquotedblleft{}As if\textquotedblright{}.) Say it once more.
\end{tcolorbox}
\section[(dialogue)]{Looking back: from How Things Work to As Long As}

\label{lesson:MR-R381-review-374-381}



\begin{quote}
Say the Marathi for a process and a phase.
\end{quote}

\subsection*{Your turn}
Say these aloud:

\begin{itemize}
\raggedright
  \item a process, a phase, a formula, a definition, a science
  \item intellect, memory power, concentration, to comprehend, to practise
  \item to weigh up, to assess, to evaluate, quality, a standard
  \item an influence, outside pressure, to encourage, inspiration, to bother
  \item a condition, an exception, a curb, a criterion, to fulfil
  \item a limit, minimum, maximum, to relax, to be manageable
  \item compulsory, optional, unavoidable, to impose, a provision
  \item as long as, until then, provided that, so that, as if
\end{itemize}

\begin{tcolorbox}[breakable,colback=teal!4,colframe=teal!35!black,title={Before you move on}]
A process? (\textbf{\mr{प्रक्रिया}}, \emph{prakriyā}.) A phase? (\textbf{\mr{टप्पा}}, \emph{ṭappā}.) A formula? (\textbf{\mr{सूत्र}}, \emph{sūtra}.) A definition? (\textbf{\mr{व्याख्या}}, \emph{vyākhyā}.) A science? (\textbf{\mr{शास्त्र}}, \emph{śāstra}.) Intellect? (\textbf{\mr{बुद्धी}}, \emph{buddhī}.) Memory power? (\textbf{\mr{स्मरणशक्ती}}, \emph{smaraṇśaktī}.) Concentration? (\textbf{\mr{एकाग्रता}}, \emph{ekāgratā}.) To comprehend? (\textbf{\mr{आकलन} \mr{होणे}}, \emph{ākalan hoṇe}.) To practise? (\textbf{\mr{सराव} \mr{करणे}}, \emph{sarāv karṇe}.) To weigh up? (\textbf{\mr{तोलणे}}, \emph{tolṇe}.) To assess? (\textbf{\mr{पारखणे}}, \emph{pārakhṇe}.) To evaluate? (\textbf{\mr{मूल्यमापन} \mr{करणे}}, \emph{mūlyamāpan karṇe}.) Quality? (\textbf{\mr{गुणवत्ता}}, \emph{guṇvattā}.) A standard? (\textbf{\mr{दर्जा}}, \emph{darjā}.) An influence? (\textbf{\mr{प्रभाव}}, \emph{prabhāv}.) Outside pressure? (\textbf{\mr{दबाव}}, \emph{dabāv}.) To encourage? (\textbf{\mr{प्रोत्साहन} \mr{देणे}}, \emph{protsāhan deṇe}.) Inspiration? (\textbf{\mr{प्रेरणा}}, \emph{prerṇā}.) To bother? (\textbf{\mr{खटकणे}}, \emph{khaṭakṇe}.) A condition? (\textbf{\mr{अट}}, \emph{aṭ}.) An exception? (\textbf{\mr{अपवाद}}, \emph{apvād}.) A curb? (\textbf{\mr{निर्बंध}}, \emph{nirbandh}.) A criterion? (\textbf{\mr{निकष}}, \emph{nikaṣ}.) To fulfil? (\textbf{\mr{पूर्तता} \mr{करणे}}, \emph{pūrtatā karṇe}.) A limit? (\textbf{\mr{मर्यादा}}, \emph{maryādā}.) Minimum? (\textbf{\mr{किमान}}, \emph{kimān}.) Maximum? (\textbf{\mr{कमाल}}, \emph{kamāl}.) To relax? (\textbf{\mr{शिथिल} \mr{करणे}}, \emph{śithil karṇe}.) To be manageable? (\textbf{\mr{झेपणे}}, \emph{jhepṇe}.) Compulsory? (\textbf{\mr{अनिवार्य}}, \emph{anivārya}.) Optional? (\textbf{\mr{ऐच्छिक}}, \emph{aicchik}.) Unavoidable? (\textbf{\mr{अपरिहार्य}}, \emph{aparihārya}.) To impose? (\textbf{\mr{लादणे}}, \emph{lādṇe}.) A provision? (\textbf{\mr{तरतूद}}, \emph{tartūd}.) As long as? (\textbf{\mr{जोपर्यंत}}, \emph{joparyant}.) Until then? (\textbf{\mr{तोपर्यंत}}, \emph{toparyant}.) Provided that? (\textbf{\mr{बशर्ते}}, \emph{baśarte}.) So that? (\textbf{\mr{जेणेकरून}}, \emph{jeṇekarūn}.) As if? (\textbf{\mr{जणू}}, \emph{jaṇū}.)
\end{tcolorbox}
